%% parts/p02/second-and-higher-order-logic-supplement.tex
% Necessity: the exact source restores full and Henkin semantics and its
% semantic translation, but it deliberately introduces no local derivability
% relation or matched calculus; the translation-defined sheet is distinct.
% Definition: this OPEN addition is exactly historical lines 338--436 followed
% by the consumer-corrected derivative of lines 458--496.
% Standing: it is editorial apparatus and a metalinguistically defined
% calculus addition, not inherited source or a full-semantics completeness claim.
% Nonredundancy: the source supplies the semantic equivalence only; this block
% adds an eliminable derivability relation and a comparison table.
% Use: the Part II driver routes it after all eleven exact higher-order
% wrappers, and the higher-order route regression pins both bounded blocks.

\noindent\textit{Reader note.}  The exact inherited chapter above fixes the
relational second-order language, full and Henkin semantics, and the basic
many-sorted semantic translation.  The following separately authored material
adds the exact schema-closed Henkin class and incidence translation required by
the translation-defined calculus.  It does not replace or alter the inherited
semantic presentation \cite{Henkin1950,Manzano1996}.

\begin{hbdefinition}{def:p02-sol-relation-substitution}{Capture-avoiding relation substitution}
For each \(n\geq1\), fix once and for all an enumeration
\(\mathsf{RelVar}_n=\{X^n_0,X^n_1,\ldots\}\); ``least fresh'' below means
least index in this enumeration absent from the displayed source formula,
inserted symbol, and substitution result constructed so far.
For an \(n\)-ary relation variable \(X^n\) and an \(n\)-ary relation variable
or predicate constant \(U^n\), the formula
\(\varphi[U^n/X^n]\) is obtained recursively by replacing each free atom
\(X^n(t_1,\ldots,t_n)\) with \(U^n(t_1,\ldots,t_n)\), commuting with Boolean
connectives and object quantifiers, and stopping beneath a binder for
\(X^n\).  Before descending beneath a binder for \(U^n\) when \(U^n\) is a
variable, rename that binder and its bound occurrences to the least fresh
relation variable of the same arity; other relation binders are traversed
recursively.  The substitution is \emph{free for} \(X^n\) precisely when this
capture-avoiding recursion is defined.  Predicate constants have no binding
occurrences and require no renaming.  The operation is metalinguistic,
preserves arity, and supplies exactly the four relation-quantifier rule
instances on the linked calculus sheet.
\index{substitution!relation variable}
\end{hbdefinition}

\section{Matched Henkin semantics and incidence translation}
\label{sec:p02-sol-matched-calculus}

Merely selecting nonempty relation domains is too weak for a calculus that
contains comprehension or choice.  The next definition fixes the exact model
class needed by the displayed sheet; the adjective ``matched'' records the
identity of its semantic closure conditions with the eligible syntactic schema
instances and has no broader meaning.

\begin{hbdefinition}{def:p02-sol-matched-henkin}{Matched Henkin models}
A relational Henkin structure as fixed in
Section~\ref{sec:second-order-logic-05} is a \emph{matched
comprehension--\allowbreak choice Henkin model}, abbreviated an
\(\mathsf{H}_{\mathsf{CC}}\)-model, iff every \(n\)-ary predicate constant
\(P\) satisfies \(P^{\mathcal H}\in\mathcal R_n\) and the structure satisfies
every eligible instance of both schema families
\[
\begin{aligned}
\mathsf{Comp}_{\chi}^{n}&\defeq
 \operatorname{ucl}\!\left(\exists X^n\,\forall\bar x
   (X^n(\bar x)\leftrightarrow\chi(\bar x,\bar p,\bar Q))\right),\\
\mathsf{Choice}_{\chi}^{n}&\defeq
 \operatorname{ucl}\!\Bigl(\Bigl(\forall\bar x\,\exists y\,
 \chi(\bar x,y,\bar p,\bar Q)\Bigr)\to{}\\
&\hspace{18mm}\exists F^{n+1}\,\forall\bar x\,\exists y\,
 \Bigl(F^{n+1}(\bar x,y)\land\chi(\bar x,y,\bar p,\bar Q)\\
&\hspace{38mm}\land
 \forall z(F^{n+1}(\bar x,z)\to z=y)\Bigr)\Bigr).
\end{aligned}
\]
Here \(n\geq1\); \(\bar x=(x_1,\ldots,x_n)\) is pairwise distinct; in
comprehension \(X^n\) is not free in \(\chi\); in choice \(F^{n+1}\) is fresh
and not free in \(\chi\); and \(\operatorname{ucl}\) universally closes
exactly the displayed object and relation parameters \(\bar p,\bar Q\).
Comprehension may be impredicative.  Choice asserts an admissible single-valued
selecting graph and introduces no function variable.  The two families are
nonredundant: comprehension collects a formula-defined relation, while choice
selects one witness from each nonempty fibre.  Their exact translations form
the background theory used below.
\index{comprehension!second-order}
\index{choice!relational graph schema}
\end{hbdefinition}

Under full semantics, eligible comprehension is automatic; eligible choice is
valid only relative to enough metatheoretic choice for the displayed families.
For matched Henkin semantics both families hold by definition, as a restriction
on the model class rather than a consequence of
\(\mathcal R_n\subseteq\mathcal P(M^n)\).

\begin{hbdefinition}{def:p02-sol-henkin-consequence}{Matched Henkin consequence}
Write \(\Gamma\entails_{\mathrm H}^{\mathsf{CC}}\varphi\) iff every
\(\mathsf H_{\mathsf{CC}}\)-model under every matching object-and-relation
assignment satisfying \(\Gamma\) also satisfies \(\varphi\).  The superscript
records comprehension--choice closure and distinguishes this relation from
consequence over arbitrary relational Henkin structures.  It is the semantic
relation paired with the derived calculus below.
\index{semantic consequence!Henkin second-order!matched schemas}
\end{hbdefinition}

The calculus registry displays its four relation-quantifier rules and two
schema families on \hyperref[sheet:sol]{the translation-defined relational
second-order sheet}.  It inherits finite-multiset sequents and object equality
from \hyperref[sheet:lk-eq]{classical first-order LK with equality}.

\begin{hbdefinition}{def:p02-sol-incidence-translation}{Incidence translation}
Assume that the nonlogical second-order signature is at most countable, so the
translated sorts and symbols use the countable variable-stock injection fixed
in Definition~\ref{def:p02-many-sorted-signature}.  For the language fixed in
Section~\ref{sec:second-order-logic-01}, form a strict
many-sorted first-order language with object sort \(D\), relation sort \(S_n\)
for each \(n\geq1\), incidence predicate \(E_n:S_n\times D^n\), and a
relation-sort constant \(\widehat P:S_n\) for each predicate constant \(P^n\).
Define \(\varphi\mapsto\varphi^\sharp\) recursively by
\[
\begin{aligned}
 (X^n(\bar t))^\sharp&\defeq E_n(X^n,\bar t),
 & (P^n(\bar t))^\sharp&\defeq E_n(\widehat P,\bar t),\\
 (\forall x\,\psi)^\sharp&\defeq\forall x{:}D\,\psi^\sharp,
 & (\exists x\,\psi)^\sharp&\defeq\exists x{:}D\,\psi^\sharp,\\
 (\forall X^n\,\psi)^\sharp&\defeq\forall X^n{:}S_n\,\psi^\sharp,
 & (\exists X^n\,\psi)^\sharp&\defeq\exists X^n{:}S_n\,\psi^\sharp.
\end{aligned}
\]
The map leaves object terms and object equality unchanged and commutes with
the classical connectives.  For a formula set \(\Gamma\), put
\(\Gamma^\sharp\defeq\{\psi^\sharp:\psi\in\Gamma\}\).  This is an eliminable
metalinguistic translation, not an object-language connective; unlike erasing
second-order binders, it makes relation application into incidence, represents
predicate constants by relation-sort objects, and keeps object and relation
quantifier ranges in different sorts.  Those distinctions are used in the
correspondence theorem.
\index{incidence translation}
\index{second-order logic!many-sorted translation}
\end{hbdefinition}

Given \(\mathcal H=(M,(\mathcal R_n)_{n\geq1},\ldots)\) such that
\(P^{\mathcal H}\in\mathcal R_n\) for every \(n\)-ary predicate constant,
define the strict many-sorted structure \(\mathcal H^\sharp\) by tagging every
sort:
\[
 D^{\mathcal H^\sharp}=\{D\}\times M,
 \qquad S_n^{\mathcal H^\sharp}=\{S_n\}\times\mathcal R_n,
\]
\[
\begin{aligned}
 E_n^{\mathcal H^\sharp}
   \bigl((S_n,R),(D,a_1),\ldots,(D,a_n)\bigr)
 &\Longleftrightarrow(a_1,\ldots,a_n)\in R,\\
 \widehat P^{\mathcal H^\sharp}&=(S_n,P^{\mathcal H}).
\end{aligned}
\]
Object constants and functions are transported along the object tag, and a
matching assignment is sent by
\(s^\sharp(x)=(D,s(x))\) and
\(s^\sharp(X^n)=(S_n,s(X^n))\).  The tags make the sorts disjoint even when
the underlying sets \(M,\mathcal R_1,\mathcal R_2,\ldots\) overlap.

\begin{hbtheorem}{thm:p02-sol-incidence-satisfaction}{proved=annex:p2.hol.7}{Incidence translation preserves satisfaction}
For every relational Henkin structure \(\mathcal H\) whose predicate-constant
interpretations belong to their declared relation domains, every relational
second-order formula \(\varphi\), and every matching assignment \(s\),
\[
 \mathcal H,s\entails\varphi
 \quad\Longleftrightarrow\quad
 \mathcal H^\sharp,s^\sharp\entails\varphi^\sharp.
\]
Moreover, \(\mathcal H\) is an \(\mathsf H_{\mathsf{CC}}\)-model iff
\(\mathcal H^\sharp\) satisfies the many-sorted translations of exactly the
eligible \(\mathsf{Comp}_{\chi}^{n}\) and
\(\mathsf{Choice}_{\chi}^{n}\) instances.
\end{hbtheorem}

% BEGIN HISTORICAL SLICE 338-436
\begin{hbdefinition}{def:p02-sol-derived-calculus}{Translation-defined SOL calculus}
Work over an at-most-countable nonlogical second-order signature and the
corresponding countable incidence signature just fixed.  Let
\(T_{\mathsf{CC}}^\sharp\) be the many-sorted first-order theory
consisting of the incidence translations of exactly the eligible
comprehension and choice instances above.  For a many-sorted formula set
\(\Delta\) and formula \(\psi\), define
\[
 \Delta\derives_{\mathsf{MSFOL}}\psi
 \quad\Longleftrightarrow\quad
 \mathsf{Ax}_{\mathrm{sort}}\cup
 \mathsf{VarGuard}_{\mathrm{MS}}(\Lang_{\mathrm{MS}})\cup
 \operatorname{tr}_{\mathrm{MS}}[\Delta]
 \derives_{\mathsf{ND}_{\mathsf{FOL}}}
 \operatorname{tr}_{\mathrm{MS}}(\psi),
\]
using Definitions~\ref{def:p02-many-sorted-translation} and the registered
first-order natural-deduction calculus.  This is again an eliminable
metalinguistic abbreviation; Theorem~\ref{thm:p02-many-sorted-translation},
its formula induction, and classical first-order adequacy make it sound and
strongly complete under the same set-theoretic hypotheses.  Now define
\[
 \Gamma\derives_{\mathsf{SOL}_{\mathrm H}^{\mathsf{CC}}}\varphi
 \quad\Longleftrightarrow\quad
 T_{\mathsf{CC}}^\sharp\cup\Gamma^\sharp
 \derives_{\mathsf{MSFOL}}\varphi^\sharp.
\]
The left side is thus an eliminable metalinguistic abbreviation for the right
side, not a newly postulated native derivation relation.  This definition is
needed because Henkin's typed completeness theorem does not cover the new
relational rule surface; it is used immediately to state an exact adequacy
theorem.  The formulas displayed on
\hyperref[sheet:sol]{the relational second-order calculus sheet} are admissible
projections of this definition, rather than an independently complete finite
list of native generators.

The advertised two-sided sheet has the following exact derived meaning.  For
finite-multiset contexts \(\Sigma,\Pi\), let \(\Sigma^\sharp\) be the
translated support of \(\Sigma\), and let \(\bigvee\Pi^\sharp\) be the
disjunction of any listing of \(\Pi^\sharp\), with
\(\bigvee\varnothing\defeq\bottom\).  Declare
\[
 \Sigma\Rightarrow\Pi\text{ derivable on the linked SOL sheet}
 \quad\Longleftrightarrow\quad
 T_{\mathsf{CC}}^\sharp\cup\Sigma^\sharp
 \derives_{\mathsf{MSFOL}}\bigvee\Pi^\sharp.
\]
Classical first-order derivability makes this independent of listing,
bracketing, and repeated occurrences.  Thus it is an eliminable sequent
abbreviation, not a second calculus.  Its assignment guard includes every
sorted object and relation variable, hence in particular every variable free
in either context.
\index{second-order calculus!translation-defined}
\end{hbdefinition}

\begin{hbtheorem}{thm:p02-sol-derived-adequacy}{proved=annex:p2.hol.8}{Strong adequacy of the matched relational calculus}
Assume the Axiom of Choice in the ambient set theory.  For every set
\(\Gamma\) of formulas and every formula \(\varphi\) in the at-most-countable
relational second-order signature fixed above,
\[
 \Gamma\derives_{\mathsf{SOL}_{\mathrm H}^{\mathsf{CC}}}\varphi
 \quad\Longleftrightarrow\quad
 \Gamma\entails_{\mathrm H}^{\mathsf{CC}}\varphi.
\]
All six additions on \hyperref[sheet:sol]{\texttt{sheet:sol}} are admissible
for this derivability relation and for its two-sided projection just defined.
The conclusion is
strong completeness for the exact matched
selected-domain semantics; it is not completeness for full validity.
\end{hbtheorem}

\begin{proofboundary}{def:p02-sol-native-boundary}{What has and has not been proved about the displayed sheet}
The theorem defines derivability by the exact many-sorted reduction and then
proves the displayed rules admissible.  It does not assert that the inherited
first-order rules, four relation-quantifier rules, and two schema displays,
taken as an independently generated native calculus, are complete.  Such a
claim would additionally require a converse deductive-correspondence theorem.
Nor does \cite{Henkin1950} supply that missing theorem: it proves soundness and
completeness for Henkin's own typed system and its general models.
\end{proofboundary}

% END HISTORICAL SLICE 338-436

% BEGIN HISTORICAL DERIVATIVE 458-496
\section{Full versus Henkin second-order semantics}
\label{sec:p02-sol-comparison}

\begin{landscape}
\begingroup\scriptsize\setlength{\tabcolsep}{4pt}
\begin{longtable}{@{}p{1.35in}p{2.95in}p{2.95in}@{}}
\caption{Full versus matched Henkin relational semantics.}
\label{tab:p02-full-vs-henkin-sol}\\
\toprule
Property & Full relational semantics & Matched Henkin relational semantics \\ \midrule
\endfirsthead
\toprule
Property & Full relational semantics & Matched Henkin relational semantics \\ \midrule
\endhead
Relation-quantifier range & Every member of \(\mathcal P(M^n)\). & A selected nonempty \(\mathcal R_n\subseteq\mathcal P(M^n)\). \\
\addlinespace
Comprehension & Automatically valid for every eligible impredicative instance. & Must be imposed as satisfaction of exactly the eligible instances, equivalently the matching definability closure. \\
\addlinespace
Choice & Valid only relative to adequate metatheoretic choice. & Must be imposed by the eligible graph-choice instances or equivalent selector closure. \\
\addlinespace
Completeness & No sound effective recursively axiomatized calculus is complete for full validity. & Under ambient AC, strong completeness holds for the exact translation-defined calculus paired with the exact schema-closed model class, Theorem~\ref{thm:p02-sol-derived-adequacy}. \\
\addlinespace
Compactness & Fails, Section~\ref{sec:second-order-logic-04}. & Holds for the finitary matched many-sorted presentation by first-order compactness. \\
\addlinespace
Downward L\"owenheim--Skolem & Fails in general; full categoricity of the uncountable ordered field obstructs countable models. & The first-order cardinal-qualified form transfers through the many-sorted presentation. \\
\addlinespace
Upward L\"owenheim--Skolem & Fails in general; full categoricity of \(\mathsf{PA}_2\) obstructs uncountable models. & The first-order form transfers when the object sort \(D\), equivalently the object domain \(M\), is infinite, with the usual language/cardinality qualifications. \\
\addlinespace
Categoricity & Can characterize infinite structures such as \(\mathbb N\) and \(\mathbb R\) up to isomorphism. & The same written axioms need not be categorical; nonstandard selected-domain models remain. \\
\bottomrule
\end{longtable}
\endgroup
\end{landscape}

In the Henkin completeness row, ``complete'' means complete for the exact
calculus/model-class pair; it never means complete for full second-order
validity.  The comparison synthesizes
\cite{Henkin1950,Manzano1996,Shapiro1991}.

% END HISTORICAL DERIVATIVE 458-496
